\documentclass[12pt, a4paper]{article}

\usepackage[utf8]{inputenc}
\usepackage[T1]{fontenc}
\usepackage[german]{babel}
\usepackage[top=2.5cm, bottom=2.5cm, left=3cm, right=2.5cm]{geometry}
\usepackage{setspace}
\onehalfspacing 
\usepackage{graphicx}
\usepackage{enumitem}

\title{
    \textbf{\Huge Project-Polearm} \\ [0.2em]
    \textbf{\Large Die Entwicklung eines Hack-and-Slash} \\ [0.6em]
    \large Programmierung, Architektur und Evaluation von Gegnerverhalten und Leveldesign \\ [0.4em]
    \normalsize Maturitätsarbeit
}
\author{Rodin Arslan, Zhile Zhu \\ \normalsize Klasse G23g}
\date{\today}

\begin{document}

\maketitle
\newpage

\tableofcontents
\newpage

\section{Einleitung}
\subsection{Motivation und Relevanz}
\subsection{Zielsetzung und Produktdefinition}
\subsection{Erkenntnisfrage}

\section{Theoretischer Hintergrund \& Analyseframework}
\subsection{Definition von Gamedesign-Kriterien}
\subsection{Das MDA-Analyseframework (Fragebogen)}
Im Folgenden wird der strukturierte Fragebogen auf Basis des MDA-Frameworks (Mechanics, Dynamics, Aesthetics) dargestellt. Die Evaluation basiert auf insgesamt 15 gezielten Fragen, um den Prototyp sowie die Referenzspiele messbar zu analysieren.

\subsubsection{Mechanics (Spielmechaniken)}
\begin{enumerate}[label=\textbf{M\arabic*}, leftmargin=1.5cm, itemsep=0.5em]
    \item \textit{[Frage 1: Z.B. Wie intuitiv ist das grundlegende Angriffssystem?]}
    \item \textit{[Frage 2: Z.B. Wie präzise fühlt sich das Treffer- und Schadensfeedback an?]}
    \item \textit{[Frage 3: Z.B. Wie unterscheidet sich das Verhalten der zwei Gegnertypen?]}
    \item \textit{[Frage 4: Z.B. Hier eigene Frage zu den Mechaniken einfügen]}
    \item \textit{[Frage 5: Z.B. Hier eigene Frage zu den Mechaniken einfügen]}
\end{enumerate}

\subsubsection{Dynamics (Spieldynamik und Interaktion)}
\begin{enumerate}[label=\textbf{D\arabic*}, leftmargin=1.5cm, itemsep=0.5em]
    \item \textit{[Frage 6: Z.B. Entsteht durch die Gegnerplatzierung ein flüssiger Spielfluss (Flow)?]}
    \item \textit{[Frage 7: Z.B. Wie variiert der Schwierigkeitsgrad im Verlauf des Kampfabschnitts?]}
    \item \textit{[Frage 8: Z.B. Bietet das Spiel genügend taktische Optionen, um Angriffen auszuweichen?]}
    \item \textit{[Frage 9: Z.B. Hier eigene Frage zu den Dynamiken einfügen]}
    \item \textit{[Frage 10: Z.B. Hier eigene Frage zu den Dynamiken einfügen]}
\end{enumerate}

\subsubsection{Aesthetics (Ästhetik und emotionales Erleben)}
\begin{enumerate}[label=\textbf{A\arabic*}, leftmargin=1.5cm, itemsep=0.5em]
    \item \textit{[Frage 11: Z.B. Wird das Kampfgefühl als befriedigend und ''stylish`` empfunden?]}
    \item \textit{[Frage 12: Z.B. Wie gut vermittelt das visuelle Feedback (z.B. Effekte) die Wucht der Treffer?]}
    \item \textit{[Frage 13: Z.B. Erinnert das Gesamtgefühl an die klassische Ära der Stylish-Action-Games?]}
    \item \textit{[Frage 14: Z.B. Hier eigene Frage zur Ästhetik einfügen]}
    \item \textit{[Frage 15: Z.B. Hier eigene Frage zur Ästhetik einfügen]}
\end{enumerate}

\subsection{Analyse der Referenzspiele}
\subsubsection{Referenzspiel 1}
\subsubsection{Referenzspiel 2}
\subsection{Ableitung der Designziele}

\section{Softwarearchitektur und technische Entwicklung}
\subsection{Systemarchitektur in Godot}
\subsection{Zentrale Gameplay-Systeme}
\subsubsection{Spielerbewegung und Kamera}
\subsubsection{Kampfsystem, Treffererkennung und Schadenslogik}
\subsection{Gegner-KI und Verhaltenslogik}

\section{Level- und Gegnerdesign}
\subsection{Gestaltung des Kampfabschnitts}
\subsection{Gegnerplatzierung (Enemy Placement)}

\section{Evaluation und Playtesting}
\subsection{Methodik des Playtests}
\subsection{Auswertung der Ergebnisse}
\subsection{Iteration aufgrund von Feedback}

\section{Schlussfolgerung \& Diskussion}
\subsection{Beantwortung der Erkenntnisfrage}
\subsection{Kritische Würdigung \& Ausblick}

\section{Reflexion}
\subsection{Reflexion des Arbeitsprozesses und der Gruppenarbeit}
\subsection{Reflexion über den Einsatz von KI-Hilfsmitteln}

\newpage
\section{Verzeichnisse}
\subsection{Literatur- und Quellenverzeichnis}
\subsection{Assetliste (Modelle, Sounds, Texturen)}
\subsection{KI-Deklaration}

\end{document}o